%% 5_real
\section{Evaluation on Real Measurements}
\label{sec:real}

\subsection{Measurement Corpus}

The measurement corpus is large but only partly usable, because evaluating
reconstruction requires repeat visits to the same place and two of its four buildings
never received them. The corpus spans four buildings, $61$ floors, $204$
transmitters, six to seven handset models per site and $5.06$ million logged
readings. Each reading carries a handset
identifier, a reference-point index, a transmitter identifier and a received
strength. Two of the four buildings, both residential towers, yield no cell with
enough co-located repeat measurements to evaluate on, so all numbers below come
from the remaining two, a university building with interior partitions and an open
multi-storey car park. The car park serves throughout as a negative control,
because an open concrete structure has the geometry of the problem without the
walls.

Dense reference points create a shortcut that scoring has to exclude, and the
thresholds we use to exclude it give rise to the two cell sets reported below. We
hold out one point and predict it from the rest of the same cell behind a $1$~m
exclusion buffer, so a held-out point is never recovered from a neighbour standing
beside it. Two cell sets appear below and
we keep them separate throughout. The reconstruction comparison uses the $54$ cells
with at least ten reference points, giving $1924$ held-out predictions. The
transmitter-locatability and residual-model analyses use the looser threshold of
eight reference points, which admits $67$ cells. Numbers from the two sets are never
compared against each other.

\begin{figure}[t]
\centering
\includegraphics[width=\columnwidth]{fig/walls.pdf}
\caption{\textbf{Walls attenuate at the treatment site and not at the control
site.} Excess absolute difference in received strength between reference-point
pairs at matched separation, split by how many walls separate them. Error bars are
bootstrap $95\%$ intervals. The open car park has the geometry of the problem
without the partitions and shows no effect.}
\label{fig:walls}
\end{figure}

\subsection{Wall Attenuation and Negative Control}

The premise behind \cref{eq:occ} can be tested on the measurements alone, without any
model, by comparing reference-point pairs at matched separation that differ in how
much wall lies between them. For pairs of reference points at matched separation we compare the excess
absolute difference in received strength between pairs separated by many walls and
pairs separated by few. \Cref{fig:walls} reports the result.

The effect appears at the walled site in every separation band wide enough to contain
a partition, and at the open control site it does not appear at all. Pairs separated by walls
differ by an additional $2.31$~dB at two to three metres, $2.50$~dB at three to
four metres and $1.51$~dB at four to six metres, with disjoint bootstrap intervals
in all three bands. At one to two metres the gap is $0.55$~dB and the intervals
overlap, which is the expected behaviour when the separation is shorter than the
spacing between partitions. At the open control site the same measurement returns
$-0.02$~dB and $+0.91$~dB with overlapping intervals and wall correlations of
$-0.057$ and $-0.017$, which is to say nothing.

The same contrast survives when the effect is put to work, and the exact null at the
control site is the strongest causal evidence in this paper that \cref{eq:occ}
identifies the right mechanism. Cross-validated
wall-aware interpolation improves accuracy by $3.86\%$ at the walled site, with the
penalty weight selected as $32$ and $16$ on the two folds. At the control site the
gain is $0.00\%$ and cross-validation selects a penalty of exactly zero on both
folds, which is the procedure declining to use geometry when there is none. A real effect where
walls exist and an exact null where they do not is what separates a causal reading
from a correlational one.

\subsection{Transmitter Locatability}
\label{sec:locatable}

The transmitter position that \cref{eq:occ} requires is recoverable in only $12$ of
the $67$ measured cells, and the corpus holds no ground truth against which even
those could be checked. Fitting a point source per cell and gating on fit quality, stability
and cross-split agreement, only $12$ of $67$ cells pass, which is $17.9\%$. The
remainder are consistent with distributed or repeated sources, which is what indoor
boosters sharing a cell identifier produce. The median coefficient of determination
across all cells is $0.31$ against a gate of $0.30$, so the count is sensitive to
where the gate is placed, and no ground-truth transmitter position exists anywhere
in the corpus against which a localisation error could be quoted.

That scarcity is the crux of the transfer problem, because the input carrying the
simulated result is available one time in five and \cref{sec:txfree_result} showed
the obvious substitute for it failing.

\begin{table}[t]
\centering\footnotesize
\begin{tabular}{@{}lrl@{}}
\toprule
Method & RMSE (dB)$\downarrow$ & vs.\ learned \\
\midrule
RBF (multiquadric) & \textbf{5.50}\,\tiny[5.26,5.77] & $p{=}$2e-30 \\
wall-aware IDW & 5.61\,\tiny[5.37,5.88] & $p{=}$1e-19 \\
IDW ($p{=}2$) & 5.79\,\tiny[5.54,6.06] & $p{=}$4e-18 \\
GP / Kriging & 5.95\,\tiny[5.69,6.22] & $p{=}$7e-07 \\
learned, with geometry & 6.20\,\tiny[5.95,6.44] & -- \\
learned, no geometry & 6.30\,\tiny[6.05,6.55] & $p{=}$7e-06 \\
\bottomrule
\end{tabular}
\caption{\textbf{Reconstruction on real indoor cellular measurements}, held-out
reference points with a $1$~m exclusion buffer, $1924$ points over $54$ cells. The
right column is a paired Wilcoxon test against the learned model with geometry.
Classical interpolation wins here, and \cref{sec:real} explains why the two
benchmarks disagree.}
\label{tab:real}
\end{table}

\subsection{Comparison with Classical Interpolation}

Classical interpolation beats the learned reconstructor on real measurements by a
margin that further training would not reverse, as \cref{tab:real} reports under
leave-one-reference-point-out cross-validation. Multiquadric radial basis interpolation reaches
$5.51$~dB against $6.20$~dB for the learned model with geometry, and the paired
Wilcoxon test over $1924$ held-out points gives $p = 2.5 \times 10^{-30}$. This is
not a marginal outcome that more training would reverse.

Two smaller results sharpen that outcome rather than soften it. Wall information
still improves the learned model, from $6.30$ to $6.20$~dB at
$p = 7.0 \times 10^{-6}$, so the mechanism is intact even where the model loses. A
residual formulation that corrects the classical prediction does win, at $5.14$
against $5.49$~dB for the same interpolator, $p = 2.1 \times 10^{-4}$. That pair uses
the $67$-cell set rather than the $54$-cell set of \cref{tab:real}, which is why the
interpolator reads $5.49$ here and $5.51$ there, and its intervals overlap, so we
report it as a paired comparison.

Classical wins here because of what the real evaluation is able to ask, not because
of a defect in the learned model. Real reference points are dense and clustered inside each cell, so
holding one out leaves neighbours a metre or two away and asks for interpolation
between nearby samples of the same process. That is the regime in which radial
basis interpolation is close to optimal by construction. The synthetic benchmark
asks a different question, namely extrapolation into free space that no trajectory
visited, and it is there that geometry conditioning pays. The two benchmarks
disagree because they are not testing the same capability, and the field's practice
of reporting only the second has allowed that distinction to stay hidden.

\subsection{Why the Two Regimes Disagree}

Read together, the three findings account for the transfer gap without appealing to
domain shift at all. The dominant error
mode is non-line-of-sight prediction, every architecture stalls at the same place,
and the input that breaks the stall is a line integral anchored at the transmitter.
On simulated data that anchor is given. On real crowdsensed cellular it is absent
four times in five, and marginalising it away destroys the signal. The gap between
the two regimes is therefore not a domain shift that fine-tuning closes. It is a
missing input, and closing it requires either transmitter infrastructure that
crowdsensing does not have or a directional geometry estimator that does not yet
exist.
